\chapter{Research Hypotheses and Experimental Design}

This chapter formalizes the research questions motivating this thesis and describes the experimental protocol designed to address them. The approach follows a two-stage structure: a broad screening benchmark to map general behavior, followed by a focused deep dive to test SHAP-intensive strategies under controlled conditions.

\section{Research Questions}

The central inquiry of this thesis can be stated as:

\begin{quote}
\textit{Does SHAP-based feature selection provide practically relevant advantages over established alternatives when applied to tabular regression, considering predictive accuracy, selection stability, and computational cost jointly?}
\end{quote}

This question decomposes into several testable components:

\begin{enumerate}
    \item \textbf{Effectiveness}: Do SHAP-selected feature subsets yield lower prediction error than subsets chosen by standard methods?
    \item \textbf{Cost}: What computational overhead does SHAP-based selection incur, and does any performance gain justify this cost?
    \item \textbf{Stability}: Are SHAP-selected subsets consistent across repeated runs, or does the method introduce additional variability?
    \item \textbf{Generalization}: Do observed patterns hold across different datasets and model families, or are conclusions dataset-specific?
\end{enumerate}

\section{Hypotheses}

Based on the literature review and preliminary analysis, the following hypotheses guide the experimental work:

\subsection{H1: Competitive but Non-Dominant Performance}

In the broad benchmark across heterogeneous datasets and model families, SHAP-based feature selection will produce competitive prediction error compared to standard methods (mutual information, recursive feature elimination, tree-based importance), but will not be uniformly dominant. Different methods will excel on different dataset-model combinations, reflecting the no-free-lunch principle in feature selection.

\subsection{H2: Marginal Gains at Intermediate Retention}

In focused evaluation with XGBoost, SHAP-based recursive elimination and hybrid strategies may produce small improvements at intermediate feature retention levels (e.g., $k=0.25$ to $k=0.50$). At very low retention ($k < 0.15$), aggressive pruning may harm all methods; at full retention ($k=1.0$), selection becomes irrelevant.

\subsection{H3: Substantial Computational Overhead}

SHAP-intensive strategies, particularly those involving iterative SHAP computation (custom SHAP-RFE, hybrid pipelines), will require substantially higher selection time than native importance or simple filter methods.

\subsection{H4: Reduced Stability at Low Retention}

Selection stability will decrease as the retention ratio decreases, across all methods. However, SHAP-based methods may exhibit additional instability due to the stochastic nature of tree ensembles and the sensitivity of SHAP rankings to model refitting.

\section{Experimental Design}

The experimental protocol consists of two stages designed to balance breadth and depth.

\subsection{Stage 1: Broad Screening Benchmark}

The first stage surveys the landscape of feature selection performance across multiple datasets, models, and methods. This screening identifies general patterns and informs the selection of configurations for deeper analysis.

\textbf{Datasets}: The four datasets described in Chapter~2 (Allstate Claims, Bike Sharing, Communities and Crime, Superconductor) span a range of sizes and dimensionalities.

\textbf{Models}: Three model families are evaluated:
\begin{itemize}
    \item Ridge regression (linear baseline)
    \item ExtraTrees (bagging ensemble)
    \item XGBoost (gradient boosting)
\end{itemize}

\textbf{Feature selection strategies}:
\begin{itemize}
    \item Mutual Information (MI): filter method based on statistical dependency
    \item Recursive Feature Elimination (RFE): wrapper method with iterative removal
    \item Tree-based Importance (Tree-FI): embedded method using native importance
    \item SHAP-based Importance: ranking by mean absolute SHAP value
\end{itemize}

\textbf{Feature retention levels}: Selection is evaluated at $k \in \{0.25, 0.50, 0.75, 1.00\}$, where $k$ represents the proportion of original features retained.

\textbf{Evaluation}: Each configuration is run with cross-validation to obtain robust error estimates. Metrics include RMSE (primary), MAE, and MAPE. Selection time and total execution time are recorded.

\subsection{Stage 2: Focused Deep Dive}

The second stage narrows focus to XGBoost---the model family showing consistently strong and operationally relevant performance in Stage~1---and tests more sophisticated SHAP-based strategies.

\textbf{Datasets}: Two datasets are selected for focused analysis:
\begin{itemize}
    \item Superconductor: moderate scale, dense numeric structure
    \item Allstate Claims: large scale, high dimensional
\end{itemize}

\textbf{Selection strategies}:
\begin{itemize}
    \item \texttt{native\_fi}: XGBoost native feature importance ranking
    \item \texttt{custom\_shap\_rfe}: recursive elimination guided by SHAP scores, removing features iteratively until the target retention is reached
    \item \texttt{hybrid\_fi\_custom\_shap\_rfe}: initial pruning using native importance to reduce the feature set, followed by SHAP-guided recursive elimination on the reduced set
    \item \texttt{baseline}: no selection, all features retained ($k=1.0$)
\end{itemize}

\textbf{Feature retention levels}: A finer grid is used: $k \in \{0.05, 0.10, 0.15, 0.25, 0.50, 1.00\}$.

\textbf{Additional metrics}: Beyond prediction error and timing, Stage~2 tracks selection stability via Jaccard similarity between feature subsets across runs, and multicollinearity indicators in the selected subsets.

\subsection{Rationale for Two-Stage Design}

A direct exhaustive comparison of all methods across all datasets with fine-grained retention levels would be computationally prohibitive and produce unwieldy results. The two-stage approach serves multiple purposes:

\begin{enumerate}
    \item \textbf{Feasibility}: Stage~1 provides broad coverage with coarse granularity; Stage~2 invests computation where it matters most.
    \item \textbf{Focus}: By narrowing to XGBoost in Stage~2, method effects are isolated from model-family effects.
    \item \textbf{Relevance}: The hybrid and custom SHAP-RFE strategies are more computationally intensive; testing them on selected configurations ensures practical relevance of the comparison.
\end{enumerate}

\section{Evaluation Metrics}

Assessment is structured around three axes, reflecting the thesis objective of evaluating practical utility:

\textbf{Predictive effectiveness}: Root Mean Squared Error (RMSE) serves as the primary metric, supplemented by Mean Absolute Error (MAE) and Mean Absolute Percentage Error (MAPE). Lower values indicate better performance.

\textbf{Selection stability}: Measured by Jaccard similarity between feature subsets selected in different cross-validation folds or repeated runs. Values range from 0 (no overlap) to 1 (identical subsets). Higher stability suggests more robust selection.

\textbf{Computational cost}: Measured by selection time (time spent on feature selection) and total execution time (selection plus model training and evaluation). Costs are compared in absolute terms and as ratios between methods.

Comparisons are always interpreted as trade-offs. A method that achieves marginally lower RMSE but requires dramatically higher computation may not be practically preferable.

\section{Implementation and Reproducibility}

All experiments are implemented in Python using scikit-learn, XGBoost, and the SHAP library. The experimental pipeline is modular, with separate components for data loading, feature selection, model training, evaluation, and result persistence.

Cross-validation uses a fixed number of folds with reproducible random seeds. Results are saved as raw CSV files containing per-fold metrics, along with aggregated summaries. Figures are generated programmatically to ensure consistency.

\section{Chapter Summary}

This chapter has established the hypotheses driving the experimental work and described a two-stage protocol designed to evaluate SHAP-based feature selection systematically. The broad screening stage covers multiple datasets and models to identify general patterns, while the focused deep dive tests SHAP-intensive strategies where they are most likely to show value. Evaluation emphasizes practical trade-offs among effectiveness, stability, and cost. The next chapter presents the experimental execution and results.
